% GENERATED FILE. Edit canonical lessons, then run npm run generate:books.
% canonical-source-hash: fnv1a64:4bfee304323952dc
% canonical-lessons: LA-C360-manifestus, LA-C360-perspicuus, LA-C360-absurdus, LA-C360-stultitia, LA-C360-constantia

\chapter{Plain, Clear and Absurd}
\label{ch:la-plain-clear-and-absurd}
\hlchaptermodality{360}

\begin{chapteropening}
\textbf{By the end of this chapter:} \emph{I can say that a thing is plain or clear, call an idea absurd or foolish, and praise steadiness of mind.}

Complete the last lesson of chapter 360: I can say that a thing is plain or clear, call an idea absurd or foolish, and praise steadiness of mind.
\end{chapteropening}

\section[manifestus, manifesta, manifestum]{manifestus, manifesta, manifestum — evident, plain; caught in the act}

\label{lesson:LA-C360-manifestus}



\begin{quote}
Before the new one: say the Latin for likely, then the Latin for probable.
\end{quote}

\subsection*{You'll want to know: manifestus, manifesta, manifestum}
\textbf{manifestus, manifesta, manifestum} — \textquotedblleft{}evident, plain; caught in the act\textquotedblright{}.

\begin{quote}
Rēs manifesta est: fūr pecūniam habet.
\end{quote}

\begin{quote}
\emph{The matter is plain: the thief has the money.}
\end{quote}

\emph{manifestus, manifesta, manifestum}. A thief \emph{manifestus} is caught red-handed. English \emph{manifest}.

\begin{cousinweb}
\emph{manifestus} is built from \emph{manus} \textquotedblleft{}hand\textquotedblright{} and an old verb \emph{fendō} \textquotedblleft{}to strike\textquotedblright{}, the one in \emph{dēfendō} and \emph{offendō} (EtymDB; Lewis \& Short): a thing you could strike with your hand, so plain to see. English \emph{manifest} comes from it.
\end{cousinweb}

\begin{grammarlens}[title={What you've built}]
One more word for telling, explaining and describing at B1.
\end{grammarlens}

\subsection*{Your turn}
\begin{itemize}
\raggedright
  \item \emph{Say it:} \emph{manifestus, manifesta, manifestum}
  \item \emph{Say it:} \emph{manifestus, manifesta, manifestum}, once more
  \item \emph{Say it:} a sentence of your own with \emph{manifestus, manifesta, manifestum}
  \item \emph{Recall:} say the Latin for likely, then the Latin for probable, then say \emph{manifestus, manifesta, manifestum} again
\end{itemize}

\begin{tcolorbox}[breakable,colback=teal!4,colframe=teal!35!black,title={Before you move on}]
What does manifestus, manifesta, manifestum mean? (\textquotedblleft{}Evident, plain; caught in the act\textquotedblright{}.) Say it once more.
\end{tcolorbox}
\section[perspicuus, perspicua, perspicuum]{perspicuus, perspicua, perspicuum — clear, transparent; evident}

\label{lesson:LA-C360-perspicuus}



\begin{quote}
Before the new one: say the Latin for probable, then the Latin for evident.
\end{quote}

\subsection*{You'll want to know: perspicuus, perspicua, perspicuum}
\textbf{perspicuus, perspicua, perspicuum} — \textquotedblleft{}clear, transparent; evident\textquotedblright{}.

\begin{quote}
Haec ratiō perspicua est omnibus.
\end{quote}

\begin{quote}
\emph{This reasoning is clear to everyone.}
\end{quote}

\emph{perspicuus, perspicua, perspicuum}. Water or glass you see through, and an argument you see through to the end. English \emph{perspicuous}.

\begin{cousinweb}
\emph{perspicuus} is built on \emph{perspiciō} \textquotedblleft{}to see through\textquotedblright{} (EtymDB; Lewis \& Short), from \emph{per} \textquotedblleft{}through\textquotedblright{} and the old verb \emph{speciō} \textquotedblleft{}to look\textquotedblright{}. English \emph{perspicuous} comes from it.
\end{cousinweb}

\begin{grammarlens}[title={What you've built}]
One more word for telling, explaining and describing at B1.
\end{grammarlens}

\subsection*{Your turn}
\begin{itemize}
\raggedright
  \item \emph{Say it:} \emph{perspicuus, perspicua, perspicuum}
  \item \emph{Say it:} \emph{perspicuus, perspicua, perspicuum}, once more
  \item \emph{Say it:} a sentence of your own with \emph{perspicuus, perspicua, perspicuum}
  \item \emph{Recall:} say the Latin for probable, then the Latin for evident, then say \emph{perspicuus, perspicua, perspicuum} again
\end{itemize}

\begin{tcolorbox}[breakable,colback=teal!4,colframe=teal!35!black,title={Before you move on}]
What does perspicuus, perspicua, perspicuum mean? (\textquotedblleft{}Clear, transparent; evident\textquotedblright{}.) Say it once more.
\end{tcolorbox}
\section[absurdus, absurda, absurdum]{absurdus, absurda, absurdum — absurd, foolish; out of tune}

\label{lesson:LA-C360-absurdus}



\begin{quote}
Before the new one: say the Latin for evident, then the Latin for clear.
\end{quote}

\subsection*{You'll want to know: absurdus, absurda, absurdum}
\textbf{absurdus, absurda, absurdum} — \textquotedblleft{}absurd, foolish; out of tune\textquotedblright{}.

\begin{quote}
Absurdum est hoc dīcere.
\end{quote}

\begin{quote}
\emph{It is absurd to say this.}
\end{quote}

\emph{absurdus, absurda, absurdum}. First a voice out of tune, then anything that jars with sense. English \emph{absurd}.

\begin{cousinweb}
\emph{absurdus} is built from \emph{ab} \textquotedblleft{}away from\textquotedblright{} and \emph{surdus} \textquotedblleft{}deaf, dull-sounding\textquotedblright{} (EtymDB; Lewis \& Short): a note that sounds wrong. English \emph{absurd} comes from it.
\end{cousinweb}

\begin{grammarlens}[title={What you've built}]
One more word for telling, explaining and describing at B1.
\end{grammarlens}

\subsection*{Your turn}
\begin{itemize}
\raggedright
  \item \emph{Say it:} \emph{absurdus, absurda, absurdum}
  \item \emph{Say it:} \emph{absurdus, absurda, absurdum}, once more
  \item \emph{Say it:} a sentence of your own with \emph{absurdus, absurda, absurdum}
  \item \emph{Recall:} say the Latin for evident, then the Latin for clear, then say \emph{absurdus, absurda, absurdum} again
\end{itemize}

\begin{tcolorbox}[breakable,colback=teal!4,colframe=teal!35!black,title={Before you move on}]
What does absurdus, absurda, absurdum mean? (\textquotedblleft{}Absurd, foolish; out of tune\textquotedblright{}.) Say it once more.
\end{tcolorbox}
\section[stultitia, stultitiae]{stultitia, stultitiae — foolishness, folly}

\label{lesson:LA-C360-stultitia}



\begin{quote}
Before the new one: say the Latin for clear, then the Latin for absurd.
\end{quote}

\subsection*{You'll want to know: stultitia, stultitiae}
\textbf{stultitia, stultitiae} — \textquotedblleft{}foolishness, folly\textquotedblright{}.

\begin{quote}
Magna stultitia est hostī crēdere.
\end{quote}

\begin{quote}
\emph{It is great folly to trust an enemy.}
\end{quote}

\emph{stultitia, stultitiae} f., the noun of \emph{stultus}, foolish. Its opposite is \emph{sapientia}.

\begin{cousinweb}
\emph{stultitia} is the noun built on \emph{stultus} \textquotedblleft{}foolish\textquotedblright{} (EtymDB; Lewis \& Short), with the ending \emph{-itia} that makes a quality.
\end{cousinweb}

\begin{grammarlens}[title={What you've built}]
One more word for telling, explaining and describing at B1.
\end{grammarlens}

\subsection*{Your turn}
\begin{itemize}
\raggedright
  \item \emph{Say it:} \emph{stultitia, stultitiae}
  \item \emph{Say it:} \emph{stultitia, stultitiae}, once more
  \item \emph{Say it:} a sentence of your own with \emph{stultitia, stultitiae}
  \item \emph{Recall:} say the Latin for clear, then the Latin for absurd, then say \emph{stultitia, stultitiae} again
\end{itemize}

\begin{tcolorbox}[breakable,colback=teal!4,colframe=teal!35!black,title={Before you move on}]
What does stultitia, stultitiae mean? (\textquotedblleft{}Foolishness, folly\textquotedblright{}.) Say it once more.
\end{tcolorbox}
\section[cōnstantia, cōnstantiae]{cōnstantia, cōnstantiae — steadiness, firmness, consistency}

\label{lesson:LA-C360-constantia}



\begin{quote}
Before the new one: say the Latin for absurd, then the Latin for foolishness.
\end{quote}

\subsection*{You'll want to know: cōnstantia, cōnstantiae}
\textbf{cōnstantia, cōnstantiae} — \textquotedblleft{}steadiness, firmness, consistency\textquotedblright{}.

\begin{quote}
Cōnstantiam eius omnēs laudant.
\end{quote}

\begin{quote}
\emph{Everyone praises his steadiness.}
\end{quote}

\emph{cōnstantia, cōnstantiae} f. A person of \emph{cōnstantia} holds the same view today as yesterday. English \emph{constancy}.

\begin{cousinweb}
\emph{cōnstantia} is the noun made from the present participle of a verb built on \emph{con-} \textquotedblleft{}together\textquotedblright{} and \emph{stō} \textquotedblleft{}to stand\textquotedblright{} (Lewis \& Short): standing firm. English \emph{constancy} and \emph{constant} come from the same family.
\end{cousinweb}

\begin{grammarlens}[title={What you've built}]
One more word for telling, explaining and describing at B1.
\end{grammarlens}

\subsection*{Your turn}
\begin{itemize}
\raggedright
  \item \emph{Say it:} \emph{cōnstantia, cōnstantiae}
  \item \emph{Say it:} \emph{cōnstantia, cōnstantiae}, once more
  \item \emph{Say it:} a sentence of your own with \emph{cōnstantia, cōnstantiae}
  \item \emph{Recall:} say the Latin for absurd, then the Latin for foolishness, then say \emph{cōnstantia, cōnstantiae} again
\end{itemize}

\begin{tcolorbox}[breakable,colback=teal!4,colframe=teal!35!black,title={Before you move on}]
What does cōnstantia, cōnstantiae mean? (\textquotedblleft{}Steadiness, firmness, consistency\textquotedblright{}.) Say it once more.
\end{tcolorbox}
